\documentclass[11pt,a4paper]{article}
\usepackage{curso}
\begin{document}
\portada{Lección}{3.4}{Límites termodinámicos de velocidad}

\noindent\textbf{Unidad 3.}\quad\textbf{Objetivos:} límite de velocidad con Wasserstein; cota de Gelbrich; eficiencia $\eta\le1$;
protocolo óptimo. \textbf{Prerrequisitos:} 3.1, 3.2, 0.3.

\section{El límite}
Para una dinámica sobreamortiguada con movilidad $\mu$ que lleva $p_0\to p_\tau$ en tiempo $\tau$, la entropía total producida (en $\kB$)
cumple
\begin{resultado}
\begin{equation}\tau\ \ge\ \frac{W_2^2(p_0,p_\tau)}{\mu\kB T\,\Sigma}\qquad\Longleftrightarrow\qquad\Sigma\ \ge\ \frac{W_2^2}{\mu\kB T\,\tau}.\end{equation}
\end{resultado}
$W_2$ es la distancia de Wasserstein-2: $W_2^2=\min_\pi\E_\pi\|x_0-x_\tau\|^2$ sobre acoplamientos $\pi$ con marginales $p_0,p_\tau$
(transporte óptimo). Rápido y eficiente a la vez es imposible: a distancia fija, $\Sigma_{\min}\propto1/\tau$ \cite{whitelam}.

\section{Gaussianas: fórmula de Gelbrich}
$W_2^2(\N(m_1,\Sigma_1),\N(m_2,\Sigma_2))=\|m_1-m_2\|^2+\Tr\big(\Sigma_1+\Sigma_2-2(\Sigma_1^{1/2}\Sigma_2\Sigma_1^{1/2})^{1/2}\big)$, exacta para
gaussianas y cota inferior para distribuciones con esas medias y covarianzas. En 1D: $(m_1-m_2)^2+(s_1-s_2)^2$.

\section{Eficiencia y protocolo óptimo}
$\eta=\frac{W_2^2}{\mu\tau\kB T\Sigma}\le1$. Para una gaussiana de varianza fija trasladada con media $m(t)$ (trampa armónica, $\kappa=\mu=\kB T=1$),
$\Sigma=\int_0^\tau\dot m^2dt\ge(\Delta m)^2/\tau=W_2^2/\tau$ por Cauchy--Schwarz, con igualdad si $\dot m$ es constante: el óptimo recorre la
geodésica a velocidad constante.
\begin{ejemplo}
Arrastre lineal de la trampa a $v=1$ durante $\tau=3$: $m(t)=t-(1-e^{-t})$, $\Delta m\approx2{,}05$, $W_2^2\approx4{,}20$,
$\Sigma=\int\dot m^2=\tau-2(1-e^{-\tau})+\frac12(1-e^{-2\tau})\approx1{,}60$, $\Sigma_{\min}=1{,}40$, $\eta\approx0{,}88$. Una simulación \evS{}
confirma $\Delta m\approx2{,}05$ y varianza $\approx1$.
\end{ejemplo}

\begin{aplicacion}
Costo mínimo de cómputos rápidos (borrado, lección 1.5); diseño de protocolos (3.5); L4: $W_2^2$ cerrado y $\eta$ para un proceso gaussiano.
\end{aplicacion}

\begin{error}
\begin{itemize}[nosep]
\item $W_2$ (Wasserstein) $\neq W$ (trabajo). \quad $\Sigma$ es la producción total (incluye $\Delta s$).
\item Cota sobreamortiguada; con inercia hay versiones distintas. \quad $\eta=1$ en tiempo finito exige el protocolo óptimo; en $\tau\to0$, $\Sigma\to\infty$.
\end{itemize}
\end{error}

\begin{thebibliography}{9}\small
\bibitem[42]{whitelam} S. Whitelam, The thermodynamic freedom of a thermodynamic computer, 2026. arXiv:2608.27938.
\bibitem[2]{peliti} L. Peliti y S. Pigolotti, \emph{Stochastic Thermodynamics}, 2021.
\bibitem[11]{seifert} U. Seifert, Stochastic thermodynamics..., 2012.
\end{thebibliography}
\end{document}
